% Generated draft: requires mathematical and source review.
\documentclass[11pt]{article}
\usepackage[margin=25mm]{geometry}
\usepackage{fontspec}
\setmainfont{Latin Modern Roman}
\usepackage{amsmath,amssymb,mathtools,unicode-math}
\setmathfont{Latin Modern Math}
\usepackage{enumitem,needspace,xurl,hyperref}
\setlength{\parindent}{0pt}
\setlength{\parskip}{5pt}
\setlength{\emergencystretch}{3em}
\newcommand{\N}{\mathbb{N}}
\newcommand{\Z}{\mathbb{Z}}
\newcommand{\Q}{\mathbb{Q}}
\newcommand{\R}{\mathbb{R}}
\newcommand{\C}{\mathbb{C}}
\newcommand{\F}{\mathbb{F}}
\newcommand{\A}{\mathbb{A}}
\newcommand{\PP}{\mathbb{P}}
\newcommand{\E}{\mathbb{E}}
\newcommand{\eps}{\varepsilon}
\newcommand{\dd}{\,\mathrm{d}}
\begin{document}
\section*{UnsolvedMath ID: 1100205}
\textbf{Source problem code:} catalogue entry
\subsection{Definitions and mathematical statement}
Let $S$ be a closed surface of genus $g\geq2$, let $V=S\times S$, and let
\[
\Sigma=\{(x,x):x\in S\}\subset V.
\]
Let $p:\widetilde V\to V$ be a finite branched cover ramified nontrivially along $\Sigma$: away from $\Sigma$ it is an ordinary finite covering, and at least one ramification index over $\Sigma$ exceeds $1$. In the usual manifold branched-cover convention, the local model transverse to the branch surface is $z\mapsto z^e$ for an integer $e\geq1$. The source does not specify a more restrictive ramification profile.

\textbf{Assertion requested by the source exercise.} For every such $S$ and $\widetilde V$, there is no $C^2$ Riemannian metric $h$ on $\widetilde V$ for which
\[
K_h(\tau)\leq0\quad\hbox{for every }x\in\widetilde V
\quad\hbox{and every two-dimensional subspace }\tau\subset T_x\widetilde V.
\]
Here $K_h$ denotes sectional curvature. The source also supplies a natural piecewise hyperbolic nonpositively curved metric on $\widetilde V$. Its original wording calls this a CAT(0) metric, whereas the research summary describes it as locally CAT(0); these are different global assertions, and only the local CAT(0) description is used as geometric background here. A locally CAT(0) metric has CAT(0) neighborhoods, meaning geodesic triangles in those neighborhoods satisfy the Euclidean comparison inequality.

The record reports that Stadler proved the asserted obstruction; the exact regularity and branched-cover hypotheses of that result are not verified by this draft.
\subsection{Short English statement}
A finite cover of $S\times S$ with nontrivial branching along its diagonal can carry the supplied locally CAT(0) piecewise hyperbolic metric, but the requested conclusion is that it admits no $C^2$ Riemannian metric of nonpositive sectional curvature.
\subsection{Sources}
\begin{itemize}
\item M. Bestvina, \emph{Questions in Geometric Group Theory} (2004), Question 2.5, PDF page 7: \url{https://www.math.utah.edu/~bestvina/eprints/questions-updated.pdf}; manually transcribed, accessed 2026-07-29. The exercise is attributed to [BGS85, p. 2]; the record gives no expansion of that reference label.
\item Stadler, \url{https://arxiv.org/abs/1312.2198} (2013); \emph{Geometric and Functional Analysis} 25 (2015), 1575--1587; LMU thesis (2014). These are the supplied research-summary references; no article or thesis title is supplied.
\item Local registry \texttt{data/unsolvedmath/problems.json}, permanent ID $1100205$, source code catalogue entry, \emph{Questions in Geometric Group Theory --- Q 2.5}. The record labels the question solved and SOLVED-IN-LITERATURE, while the source-list background marks current-status and rights checks NEEDS\_REVIEW.
\end{itemize}
\end{document}
